\documentclass[10pt, letterpaper]{article}

\usepackage[
    ignoreheadfoot,
    top=0.8cm, bottom=0.8cm, left=1.2cm, right=1.2cm,
    footskip=1.0 cm,
    ]{geometry}
\usepackage{titlesec}
\usepackage{tabularx}
\usepackage{array}
\usepackage[dvipsnames]{xcolor}
\definecolor{primaryColor}{RGB}{0, 0, 0}
\usepackage{enumitem}
\usepackage{fontawesome5}
\usepackage{amsmath}
\usepackage[
    pdftitle={Andreas Tzitzikas's ASIC/FPGA Resume},
    pdfauthor={Andreas Tzitzikas},
    colorlinks=true,
    urlcolor=primaryColor
]{hyperref}
\usepackage[pscoord]{eso-pic}
\usepackage{calc}
\usepackage{bookmark}
\usepackage{lastpage}
\usepackage{changepage}
\usepackage{paracol}
\usepackage{ifthen}
\usepackage{needspace}
\usepackage{iftex}
\usepackage{datetime}

\ifPDFTeX
    \input{glyphtounicode}
    \pdfgentounicode=1
    \usepackage[T1]{fontenc}
    \usepackage[utf8]{inputenc}
    \usepackage{lmodern}
\fi

\usepackage{charter}

\raggedright
\AtBeginEnvironment{adjustwidth}{\partopsep0pt}
\pagestyle{empty}
\setcounter{secnumdepth}{0}
\setlength{\parindent}{0pt}
\setlength{\topskip}{0pt}
\setlength{\columnsep}{0.2cm}
\pagenumbering{gobble}

\titleformat{\section}{\needspace{4\baselineskip}\bfseries\large}{}{0pt}{}[\vspace{1pt}\titlerule]

\titlespacing{\section}{
    -1pt
}{
    0.2 cm
}{
    0.1 cm
}

\renewcommand\labelitemi{$\vcenter{\hbox{\small$\bullet$}}$}
\newenvironment{highlights}{
    \begin{itemize}[
        topsep=0.05 cm,
        parsep=0.05 cm,
        partopsep=0pt,
        itemsep=0pt,
        leftmargin=0 cm + 10pt
    ]
}{
    \end{itemize}
}

\newenvironment{onecolentry}{
    \begin{adjustwidth}{
        0 cm + 0.00001 cm
    }{
        0 cm + 0.00001 cm
    }
}{
    \end{adjustwidth}
}

\newenvironment{twocolentry}[2][]{
    \onecolentry
    \def\secondColumn{#2}
    \setcolumnwidth{\fill, 4.5 cm}
    \begin{paracol}{2}
}{
    \switchcolumn \raggedleft \secondColumn
    \end{paracol}
    \endonecolentry
}

\newenvironment{header}{
    \setlength{\topsep}{0pt}\par\kern\topsep\centering\linespread{1.5}
}{
    \par\kern\topsep
}

\let\hrefWithoutArrow\href

\begin{document}
    \newcommand{\AND}{\unskip
        \cleaders\copy\ANDbox\hskip\wd\ANDbox
        \ignorespaces
    }
    \newsavebox\ANDbox
    \sbox\ANDbox{$|$}

    \begin{header}
        \fontsize{15 pt}{15 pt}\selectfont Andreas Tzitzikas

        \vspace{3 pt}

        \normalsize
        \mbox{College Park, MD}
        \kern 3.0 pt
        \AND
        \kern 3.0 pt
        \mbox{U.S. Citizen}
        \kern 3.0 pt
        \AND
        \kern 3.0 pt
        \mbox{\hrefWithoutArrow{mailto:andreas.tz.work@gmail.com}{andreas.tz.work@gmail.com}}
        \kern 3.0 pt
        \AND
        \kern 3.0 pt
        \mbox{\hrefWithoutArrow{https://linkedin.com/in/andreas-tzitzikas}{linkedin.com/in/andreas-tzitzikas}}
        \kern 3.0 pt
        \AND
        \kern 3.0 pt
        \mbox{\hrefWithoutArrow{https://github.com/Wafer0}{github.com/Wafer0}}
    \end{header}

    \vspace{3 pt}

    \section{Education}

    \vspace{0.10cm}
    \begin{twocolentry}{Expected May 2027}
        \textbf{University of Maryland}, College Park, MD
    \end{twocolentry}
    \begin{twocolentry}{\textbf{GPA: 4.00}}
        \textit{Master of Science in Computer Engineering}
    \end{twocolentry}
    \vspace{0.05cm}
    \begin{twocolentry}{Expected May 2026}
        \textit{Bachelor of Science in Computer Engineering} (GPA: 3.88)
    \end{twocolentry}
    \vspace{0.05cm}
    \begin{onecolentry}
        \textbf{Relevant Coursework:} Digital Computer Design (ENEE646), Digital CMOS VLSI Design (ENEE640), Modern Digital System Design (ENEE408C), Parallel Computing (CMSC416), Math Foundations (ENEE641), Reverse Eng. \& HW Security (ENEE459B), CAD Tools (ENEE459A)
    \end{onecolentry}

    \section{Projects}

        \begin{twocolentry}{
            Nov 2025 - Dec 2025 } \textbf{Tomasulo Out-of-Order RISC-V CPU | SystemVerilog, OpenLane ASIC Flow}
        \end{twocolentry}
        \vspace{0.05 cm}
        \begin{onecolentry}
            \begin{highlights}
                \item Architected complete out-of-order CPU using \textbf{Tomasulo's Algorithm} with register renaming, 8-entry reservation stations, 16-entry reorder buffer, and common data bus for dynamic scheduling and speculative execution.
                \item Achieved \textbf{36\% higher frequency} (69.7MHz vs 51.2MHz) and 40-150\% IPC improvement on complex workloads compared to in-order implementation through comprehensive benchmarking and performance analysis.
                \item Synthesized to fabrication-ready GDSII using \textbf{Sky130 130nm PDK}, achieving 0 DRC violations, 0 LVS errors, and meeting all PPA constraints with comprehensive static timing analysis.
                \item Engineered a Bash-driven regression suite comparing architectural state against an ISA simulator (Golden Reference), achieving 100\% pass rate across 15 unit tests and 7 full-system workloads.
            \end{highlights}
        \end{onecolentry}

        \vspace{0.10 cm}
        \begin{twocolentry}{
            Aug 2025 - Nov 2025 } \textbf{5-Stage In-Order RISC-V CPU | SystemVerilog, ASIC Implementation}
        \end{twocolentry}
        \vspace{0.05 cm}
        \begin{onecolentry}
            \begin{highlights}
                \item Designed and implemented complete 5-stage pipelined CPU supporting full \textbf{RV32I ISA} with hazard detection, data forwarding, and comprehensive branch/jump support across 10 RTL modules.
                \item Achieved \textbf{100\% functional verification} through extensive testbench development and automated testing framework, validating all instruction types and pipeline corner cases.
                \item Synthesized to \textbf{10,121 standard cells} with 1.0mm² die area using OpenLane flow, achieving 50MHz operation with zero timing violations and DRC/LVS clean layout.
                \item Optimized critical path timing through strategic pipeline balancing, synchronous reset implementation, and comprehensive static timing analysis for ASIC manufacturability.
            \end{highlights}
        \end{onecolentry}

        \vspace{0.10 cm}

        \begin{twocolentry}{
            Aug 2025 – Dec 2025
        }
            \textbf{Polynomial Evaluation Accelerator | SystemVerilog, FPGA Synthesis}
        \end{twocolentry}
        \vspace{0.05 cm}
        \begin{onecolentry}
            \begin{highlights}
                \item Architected a high-throughput LIDE (Lightweight Dataflow Environment) accelerator; designed 5 coordinating FSMs to manage data tokens between 9 actor modules, ensuring deadlock-free operation.
                \item Achieved 100\% functional coverage through comprehensive SystemVerilog testbench development with 15+ complex test cases validated against C reference models.
                \item Optimized memory access with dual-port coefficient memory and auto-padding logic, enabling batch processing of up to 31 polynomial evaluations with minimal latency.
            \end{highlights}
        \end{onecolentry}

    \section{Experience}

        \begin{twocolentry}{
            Jun 2025 – Aug 2025
        }
            \textbf{Embedded Systems Intern}, Alchemity -- College Park, MD
        \end{twocolentry}
        \vspace{0.05 cm}
        \begin{onecolentry}
            \begin{highlights}
                \item Accelerated validation of a hybrid solid oxide fuel cell reactor by developing a full-stack automation suite (\textbf{Rust}, \textbf{Python}) to parallelize testing, converting multi-day manual experiments into overnight runs to significantly increase throughput.
                \item Engineered real-time \textbf{STM32} firmware to precisely control material synthesis within modular reactors, featuring non-blocking motor/relay drivers and \textbf{SPI} peripheral management.
                \item Built end-to-end control interfaces including desktop GUIs, embedded display drivers, and CLI tools to streamline workflows between hardware and software teams.
            \end{highlights}
        \end{onecolentry}

        \vspace{0.10 cm}

        \begin{twocolentry}{
            Jan 2024 – Present
        }
            \textbf{Technical Coordinator}, Electronics Prototyping Lab, Terrapin Works -- College Park, MD
        \end{twocolentry}
        \vspace{0.05 cm}
        \begin{onecolentry}
            \begin{highlights}
                \item Provide end-to-end \textbf{PCB} support (schematic to assembly) for 20+ students and faculty per semester, maintaining \textbf{LPKF} tools to ensure \textbf{90\% uptime}.
                \item Lead training for new employees and expand campus outreach to increase awareness of PCB services, while assisting researchers with diagnostics.
            \end{highlights}
        \end{onecolentry}

    \section{Skills}
        \begin{onecolentry}
            \textbf{HDLs \& Languages:} SystemVerilog, Verilog, C/C++, Python, Tcl, Bash \\
            \textbf{Computer Architecture:} RISC-V ISA, Pipelined Datapaths, Out-of-Order Execution, Hazard Control, RTL Design \\
            \textbf{ASIC/FPGA Tools:} OpenLane, Yosys, Sky130 PDK, Vivado, Icarus Verilog, Verible, GTKWave, STA, DRC/LVS Verification \\
            \textbf{Methodologies:} PPA Optimization, Testbench Development, Functional Verification, Automated Regression, Bash Scripting, Immediate Assertions, Waveform Debugging, Git, Linux
        \end{onecolentry}

\end{document}
